% Auto-generated deferred proofs (build_submission.py). Do not edit by hand.
% Each proof is titled by the result it proves; labels resolve to the body statements.

\begin{proof}[Proof of Theorem~\ref{thm:imp}]
\emph{Witness (non-vacuous).} Let $X\sim\mathcal N(0,1)$ under \emph{both} worlds, so
$P_T^1(X)=P_T^2(X)$. Take the fixed predictors $f_0(x)=\mathbf 1[x>0]$ and
$f_a(x)=\mathbf 1[x<0]$, and $0/1$ loss. Define
$P_T^1(Y\mid X)\colon Y=\mathbf 1[X>0]$ and $P_T^2(Y\mid X)\colon Y=\mathbf 1[X<0]$.
Then $R_T^1(f_0)=0,\ R_T^1(f_a)=1$, so $\Delta^1=-1<0$ (freeze is correct); and
$R_T^2(f_0)=1,\ R_T^2(f_a)=0$, so $\Delta^2=+1>0$ (adapt is correct). The evidence
$Z=\phi(X_{1:m},f_0,f_a)$ is a function of $X$ and of the fixed maps $f_0,f_a$ only;
since $X$ has the same law in both worlds, $\Law(Z\mid P_T^1)=\Law(Z\mid P_T^2)$.
Thus (i) and (ii) hold simultaneously.

\emph{Impossibility.} A label-free rule is a measurable map of $Z$ and independent
randomness, so $\Law(g)$ coincides across the two worlds. Suppose $g$ outputs
\textsc{adapt} with probability $q$ and \textsc{freeze} with probability $r$ on this
evidence. Measuring committal regret as the excess $0/1$ risk of a committed action
versus the better committed action in that world (and assigning abstention a separate
coverage cost), world~1 charges $q\cdot|\Delta^1|=q$ (adapting is wrong) and world~2
charges $r\cdot|\Delta^2|=r$ (freezing is wrong). The worst case is $\max(q,r)$, which
is minimized over the simplex at $q=r=0$, i.e.\ \textsc{abstain}. \qedhere
\end{proof}

\begin{proof}[Proof of Theorem~\ref{thm:imp-quant}]
A committal rule on $Z$ is a randomized test of $P_{T,Z}^{1,(n)}$ vs
$P_{T,Z}^{2,(n)}$; identifying $\{g=\textsc{adapt}\}$ with deciding world~2, $M(g)$ is
the sum of its type-I and type-II errors. The Le Cam testing-affinity identity
$\inf_\psi\{\E_{Q_1}\psi+\E_{Q_2}(1-\psi)\}=1-\TV(Q_1,Q_2)$ (attained by
$\psi^\star=\mathbf1[dQ_2>dQ_1]$) with $Q_j=P_{T,Z}^{j,(n)}$ gives the equality. Since
$Z$ is a measurable map of the observables, the data-processing inequality for total
variation gives the stated upper bound by the full-data TV. In the witness, $Z$ depends
on the worlds only through $\Law(X_{1:n})$, which coincide, so the evidence TV is $0$
and $\inf_g M(g)=1$; adding a coverage cost $c\in(0,\tfrac12)$ makes \textsc{abstain}
strictly optimal. The detectable rate is the Gaussian-location TV. \qedhere
\end{proof}

\begin{proof}[Proof of Theorem~\ref{thm:gate}]
At each $z$ the conditional risk of action $g(z)\in\{0,1\}$ is affine with
$\rho(1;z)-\rho(0;z)=-\Delta(z)$, so the excess of $\widehat g(z)$ over $g^\star(z)$ is
$|\Delta(z)|\,\mathbf1\{\widehat g(z)\neq g^\star(z)\}$ (gap $|\Delta(z)|$ on a sign
disagreement, $0$ otherwise). Take $\E_Z$; all terms are nonnegative so $g^\star$ is
optimal. \qedhere
\end{proof}

\begin{proof}[Proof of Proposition~\ref{prop:gate-minimax}]
$\widehat g$ depends on the world only through $Z$, whose law is identical, so its
adapt-probability $q(z)$ is common to both worlds; the two worlds' sign-error
probabilities are $q(z)$ and $1-q(z)$, summing to $1$, so the larger is $\ge\tfrac12$.
Hence $\max_j \mathrm{reg}_j\ge\tfrac12(\mathrm{reg}_1+\mathrm{reg}_2)=\E[|\Delta(Z)|]/2$
by Theorem~\ref{thm:gate}. \qedhere
\end{proof}

\begin{proof}[Proof of Theorem~\ref{thm:cert}]
On the event $\{|\widehat\Delta-\Delta|\le\varepsilon\}$ (probability $\ge1-\alpha$),
\textsc{adapt} implies $\Delta\ge\widehat\Delta-\varepsilon>0$, so a false adapt can
occur only on the complementary event of mass $\le\alpha$. The freeze bound is
symmetric. \qedhere
\end{proof}

\begin{proof}[Proof of Corollary~\ref{cor:forced-abstain}]
Because the two evidence laws coincide, the rule's adapt- and freeze-probabilities are
common to both worlds, say $q_a,q_f$. In world~1 ($\Delta^1<0$) adapting is a false
adapt, so $q_a\le\alpha$; in world~2 ($\Delta^2>0$) freezing is a false freeze, so
$q_f\le\alpha$. Hence $\Pr[\textsc{abstain}]=1-q_a-q_f\ge1-2\alpha$. \qedhere
\end{proof}

\begin{proof}[Proof of Corollary~\ref{cor:samplecomp}]
The rule commits when $\varepsilon(n)<|\widehat\Delta|$; for $\widehat\Delta\to\Delta$
this is $\varepsilon(n)<|\Delta|$. The dominant term $\sigma\sqrt{2\log(2/\alpha)/n}<|\Delta|$
rearranges to $n>2\sigma^2\Delta^{-2}\log(2/\alpha)$; retaining the Bernstein $1/n$ term
adds the $3(b-a)|\Delta|^{-1}\log(2/\alpha)$ summand. \qedhere
\end{proof}

\begin{proof}[Proof of Theorem~\ref{thm:pos}]
Change of measure: $\E_{P_T}[\ell(f(X),Y)]=\E_{P_S}[\tfrac{p_T(X)}{p_S(X)}\ell(f(X),Y)]
=\E_S[r(X)\ell(f(X),Y)]$, using $P_S(Y\mid X)=P_T(Y\mid X)$. Unbiasedness is immediate;
by the strong law and $\E_S[r^2]<\infty$ the estimator is consistent with variance
$O(1/n)$. Applying this to $f_0,f_a$ gives a consistent $\widehat\Delta$, whose sign
matches $\sign\Delta$ for large $n$ when $|\Delta|>0$. \qedhere
\end{proof}

\begin{proof}[Proof of Theorem~\ref{thm:disagree}]
Off $D$ the predictors agree, so $\ell(f_0(X),Y)=\ell(f_a(X),Y)$ and the per-sample
benefit $\delta=0$. On $D$ the two predictions differ, so for binary $Y$ exactly one
of $f_0,f_a$ equals $Y$; hence $\delta=+1$ when $f_a$ is correct and $-1$ when $f_0$
is correct, giving $\E[\delta\mid D]=P_T(f_a{=}Y\mid D)-P_T(f_0{=}Y\mid D)=2a_a^D-1$
(using $P_T(f_0{=}Y\mid D)=1-a_a^D$). Then
$\Delta=\E_{P_T}[\delta]=P_T(D)\,\E[\delta\mid D]=P_T(D)(2a_a^D-1)$, and since
$P_T(D)\ge0$ the sign claim follows. \qedhere
\end{proof}

\begin{proof}[Proof of Theorem~\ref{thm:anytime}]
The bet cap gives $1+\lambda_t a\ge1-c>0$, so $E_t^{+}>0$. As $\lambda_t$ is predictable,
$\E[E_t^{+}\mid\mathcal F_{t-1}]=E_{t-1}^{+}(1+\lambda_t\Delta)\le E_{t-1}^{+}$ under
$\Delta\le0$; thus $(E_t^{+})$ is a nonnegative supermartingale with $\E[E_0^{+}]=1$, and
Ville's inequality gives $\Pr[\exists t: E_t^{+}\ge1/\alpha]\le\alpha$. The freeze bound
is symmetric. \qedhere
\end{proof}

\begin{proof}[Proof of Theorem~\ref{thm:disagree-mc}]
$\delta=\mathbf1[f_0\neq Y]-\mathbf1[f_a\neq Y]$; off $D$ both indicators coincide.
On $D$, $\E[\delta\mid D]=(1-p_0)-(1-p_a)=p_a-p_0$, and $\Delta=P_T(D)\E[\delta\mid D]$.
For $K=2$, exactly one predictor is correct on $D$ so $p_0=1-p_a$ and this reduces to
Theorem~\ref{thm:disagree}; for $K\ge3$ \emph{both} may err, but the identity still holds. \qedhere
\end{proof}

\begin{proof}[Proof of Theorem~\ref{thm:disagree-reg}]
Difference of squares for $\delta$; off $D$, $f_0=f_a\Rightarrow\delta=0$; linearity gives
the split and $\Delta=P_T(D)\E[\delta\mid D]$. \qedhere
\end{proof}

\begin{proof}[Proof of Lemma~\ref{lem:reduction}]
Off $D$, $f_0=f_a$ so the pointwise excess loss $\delta(x)=\ell(f_0,Y)-\ell(f_a,Y)=0$. On
$D$ (binary, $0/1$ loss) exactly one of $f_0,f_a$ equals $Y$, so
$\delta=\mathbf 1[f_a{=}Y]-\mathbf 1[f_0{=}Y]=2\,\mathbf 1[f_a{=}Y]-1$, giving
$\E[\delta\mid D]=2\eta_a-1$ pointwise and $\E_{\mu\mid D}[\delta]=2\bar a-1$. Thus
$\Delta=\E_T[\delta]=\mu(D)\,\E_{\mu\mid D}[\delta]=2\mu(D)(\bar a-\tfrac12)$. Linearity gives
$\bar a-\tfrac12=\E_{\mu\mid D}[\eta_a]-\tfrac12=(\E_{\mu\mid D}[s]-\tfrac12)+\E_{\mu\mid D}[\eta_a-s]=M+\gamma$.
Since $\mu(D)>0$, $\sign\Delta=\sign(M+\gamma)$.
\end{proof}

\begin{proof}[Proof of Theorem~\ref{thm:frontier}]
(i) Immediate from Lemma~\ref{lem:reduction}: $\sign\Delta$ is a function of $M+\gamma$, and
$M$ is fixed by the observables. (ii) If $|M|>\beta$ then for all $|\gamma|\le\beta$,
$\sign(M+\gamma)=\sign M$ (the perturbation cannot cross $0$), so the sign is constant on
$\mathcal C_\beta$ and equals $\sign M$. Conversely if $|M|\le\beta$, recall
$\gamma=\bar a-\tfrac12-M\in[-(M+\tfrac12),\,\tfrac12-M]$, so the admissible-and-budgeted
range is $\gamma\in[\,\max(-\beta,-(M+\tfrac12)),\ \min(\beta,\tfrac12-M)\,]$. Its upper
endpoint gives $M+\gamma=\min(M+\beta,\tfrac12)\ge0$ (as $M\ge-\beta$) and its lower
endpoint gives $M+\gamma=\max(M-\beta,-\tfrac12)\le0$ (as $M\le\beta$), so the two
admissible targets have $\Delta\ge0$ and $\Delta\le0$
respectively; the sign is not constant, hence not identifiable (at $|M|=\beta$ one
endpoint gives $\Delta=0$, consistent with the strict criterion $|M|>\beta$).
(iii) $\sign\Delta=\sign(M+\gamma)$ with $M$ fixed; constancy of $\sign(M+\gamma)$ on
$\mathcal C$ is, by definition, the statement that $M+\gamma$ does not change sign on
$\mathcal C$, i.e.\ $\gamma$ stays on one side of $-M$.
\end{proof}

\begin{proof}[Proof of Theorem~\ref{thm:minimax}]
For $|M|>\beta$, Theorem~\ref{thm:frontier}(ii) makes the sign certain, so $g\equiv\sign M$
errs with probability $0$. For $|M|\le\beta$, consider the two targets $P^+$ ($\gamma=+\beta$)
and $P^-$ ($\gamma=-\beta$) of Theorem~\ref{thm:frontier}(ii); they differ only in the
conditional law $\eta_a$ on $D$, which is \emph{not} part of the observables, so their
observable laws are identical (total variation $0$). Any $g$ is a function of the
observables and therefore returns the \emph{same} action on $P^+$ and $P^-$, which have
opposite $\sign\Delta$; thus $g$ is wrong on exactly one of them, giving worst-case error
$\ge\tfrac12$, and $\tfrac12$ is achieved by either constant rule. This is the Le Cam
two-point bound $\tfrac12(1-\mathrm{TV})$ with $\mathrm{TV}=0$.
\end{proof}

\begin{proof}[Proof of Proposition~\ref{prop:beta0}]
Each method thresholds an observable surrogate $s$ for $\eta_a$ and predicts the benefit
sign by whether the average surrogate exceeds chance, i.e.\ by $\sign(\E_{\mu\mid D}[s]-\tfrac12)=\sign M$;
this is Theorem~\ref{thm:minimax} with $\beta=0$. Correctness $\sign M=\sign(M+\gamma)$ fails
exactly when $\gamma$ has the opposite sign to $M$ and $|\gamma|>|M|$ (Lemma~\ref{lem:reduction}).
\end{proof}

\begin{proof}[Proof of Theorem~\ref{thm:caltransfer}]
$p_h=\E_T[\Pr_T(h{=}Y\mid \mathrm{conf}_h,D)\mid D]$; by the calibration-transfer bound the inner
conditional is within $\eta$ of $\mathrm{conf}_h$ pointwise, so $|p_h-q_h|\le\eta$. Adding the two
$\eta$-balls gives the bracket; if its half-width $2\eta$ does not straddle $0$ the sign of
$q_a-q_0$ equals $\sign(p_a-p_0)=\sign\Delta$ (Theorem~\ref{thm:disagree-mc}). \qedhere
\end{proof}

\begin{proof}[Proof of Theorem~\ref{thm:router-rad}]
Symmetrization with the bounded-difference (McDiarmid) inequality gives the uniform deviation
bound; a depth-$d$ tree on $p$ features realizes $O(d\log p)$ effective decisions, so a $T$-term
ensemble has Rademacher complexity $O(\sqrt{Td\log p/n})$. Setting the deviation $\le\varepsilon$
(resp.\ $\le|\Delta|$) yields the radius-domination (resp.\ sample-complexity) claim. \qedhere
\end{proof}

\begin{proof}[Proof of Theorem~\ref{thm:anytime-drift}]
Under $\Delta\le0$, $\E[1+\lambda_tX_t\mid\mathcal F_{t-1}]=1+\lambda_t\mu_t\le1+\lambda_t\beta$, so
$\widetilde E_t^{+}$ has $\E[\widetilde E_t^{+}\mid\mathcal F_{t-1}]\le\widetilde E_{t-1}^{+}$ and unit
initial value; Ville's inequality gives $\Pr[\exists t:\widetilde E_t^{+}\ge1/\alpha]\le\alpha$,
i.e.\ the inflated threshold on $E_t^{+}$. \qedhere
\end{proof}

\begin{proof}[Proof of Theorem~\ref{thm:reg-noise}]
$R_T(f)=\E(f-g-\epsilon)^2=\E(f-g)^2+\E[\epsilon^2]$, the cross term vanishing by
$\E[\epsilon\mid X]=0$. The noise term cancels in the difference and survives in the
levels.
\end{proof}

\begin{proof}[Proof of Theorem~\ref{thm:reg-iff}]
Expanding the squared losses, $\delta=(f_0-f_a)(f_0+f_a-2Y)$ and taking target
expectations with $Y=g_S+b+\epsilon$ (noise vanishing by
Theorem~\ref{thm:reg-noise}) gives the displayed decomposition. H\"older yields
$|\E_T[(f_0-f_a)b]|\le \|b\|_\infty\,\E_T|f_0-f_a|\le BW$, with equality iff
$b=-B\sign(f_0-f_a)$ a.e.\ on $\{f_0\ne f_a\}$ --- an admissible member of the family,
so the bracket is tight. (i): on the $1-\alpha$ event
$|\widehat T_S-T_S|\le\varepsilon_n$, the committed sign matches
$\sign(U-2T_S)$ and $|U-2T_S|>2BW$, so no admissible drift can flip $\Delta$.
(ii): the map $t\mapsto\E_T[(f_0-f_a)\,t\,b^\star]=-tBW$ is continuous and onto
$[-BW,BW]$; pick $t_\pm$ so that $\Delta(t_\pm)=U-2T_S+2t_\pm BW$ has opposite signs.
Since $b$ enters only through the target labels, neither the unlabeled target evidence
nor the source sample distinguishes the two worlds; apply Theorem~\ref{thm:imp}.
\end{proof}

\begin{proof}[Proof of Theorem~\ref{thm:ls-iff}]
$\Delta$ is affine on the compact convex set $A(\pi)$, so
$\{\Delta(\pi'):\pi'\in A(\pi)\}$ is the interval
$[\Delta(\pi)-\rho,\,\Delta(\pi)+\rho]$ with the supremum attained.
(i) On the $1-\alpha$ event the estimate lands within $\varepsilon_n$ of a member of
$A(\pi)$; if $|\delta\cdot\widehat\pi|>\rho+\varepsilon_n$ then every member of
$A(\pi)$ --- in particular the truth --- has benefit of the committed sign.
(ii) If $|\Delta(\pi)|<\rho$ the interval straddles $0$; pick
$t_\pm$ with $\Delta(\pi+t_\pm v)$ of opposite signs ($t\mapsto\Delta(\pi+tv)$ is
affine, hence onto the interval). Worlds $\pi$ and $\pi+tv$ share the observable law by
$v\in\ker M$ and share the source by construction; apply Theorem~\ref{thm:imp}.
\end{proof}

\begin{proof}[Proof of Proposition~\ref{prop:reach-mono}]
$A'(\pi)=(\pi+\ker M')\cap\Delta_K\subseteq A(\pi)$, and the supremum over a subset is
no larger.
\end{proof}

\begin{proof}[Proof of Theorem~\ref{thm:unify}]
(i) On the $1-\alpha$ event, $|\widehat c-c|\le\varepsilon_n$, so
$|\widehat c|>\rho+\varepsilon_n$ implies $|c|>\rho$, and every $Q\in[P]$ has
$\Delta(Q)\in[c-\rho,c+\rho]$, an interval not containing $0$, of the committed sign.
(ii) Both signs are attained on $I(P)$ by definition of interior; pick two attaining
worlds; their observable laws coincide by membership in $[P]$, and
Theorem~\ref{thm:imp}'s argument (identical action law across the pair) forbids a valid
commitment.
\end{proof}

\begin{proof}[Proof of Theorem~\ref{thm:multicand}]
On $D$ with binary $Y$, $f_a^{(i)}{=}f_a^{(j)}$ iff both are correct or both wrong; under
Definition~\ref{def:cei},
$A_{ij}=a_ia_j+(1-a_i)(1-a_j)=\tfrac12(1+b_ib_j)$, giving $c_{ij}=b_ib_j$. The three
relations among $i,k,l$ give $c_{ik}c_{il}/c_{kl}=b_i^2$; positivity of the products under
the anchor yields $b_i$, hence $a_i=(1+b_i)/2$. Ordering and pairwise sign differences
follow.
\end{proof}

\begin{proof}[Proof of Proposition~\ref{prop:cei-test}]
Rank-one consistency of $c=bb^\top$ forces all $2\times2$ minors to vanish; the listed
equalities are these minors. Conversely $\tau>0$ contradicts rank one, so no advantage
vector reproduces the agreements: Definition~\ref{def:cei} fails.
\end{proof}

\begin{proof}[Proof of Theorem~\ref{thm:frontier}]
(i) The error bounds are Theorem~\ref{thm:cert}. On the event
$G=\{|\widehat\Delta-\Delta|\le\varepsilon_\alpha\}$, which has probability at least
$1-\alpha$, $\kappa_\alpha(z)>0$ implies
$|\widehat\Delta|\ge|\Delta|-\varepsilon_\alpha>\varepsilon_\alpha$, so the rule commits.
Hence $\{\kappa_\alpha>0\}\cap G\subseteq\{\text{commit}\}$ and
$C\ge\Pr[\kappa_\alpha>0]-\alpha$.
(ii) A label-free rule is a randomized functional of $Z$; by the data-processing
inequality for total variation its action laws across the pair differ by at most $t$.
In the world with $\Delta\ge\delta$ the wrong commit is \textsc{freeze}, so $f_1\le\alpha$;
in the world with $\Delta\le-\delta$ the wrong commit is \textsc{adapt}, so $a_2\le\alpha$.
Then $a_1\le a_2+t\le\alpha+t$, giving commit probability
$a_1+f_1\le2\alpha+t$ in world~1, and symmetrically in world~2.
\end{proof}

\begin{proof}[Proof of Proposition~\ref{prop:s1}]
With $\varepsilon\equiv0$ every instance is committed and all harmful mass is adapted
into; the regret expression is the false-adapt term of the decomposition. On the
ambiguous pair the rule commits \textsc{adapt} in both worlds and is fully wrong in the
$\Delta<0$ world, paying $\delta=2\cdot(\delta/2)$.
\end{proof}

\begin{proof}[Proof of Proposition~\ref{prop:s2}]
Construct the pair as in Theorem~\ref{thm:imp} with the confidence statistic a function
of $X$ alone, concentrated above the threshold in both worlds. The filter's action law is
then identical across worlds and commits \textsc{adapt} with probability $\approx1$, so in
the harmful world $\mathrm{FA}$ equals the harmful mass. The certificate's bound is
world-independent (Theorem~\ref{thm:cert}).
\end{proof}

\begin{proof}[Proof of Proposition~\ref{prop:s3}]
$a'=(1-\eta)a+\eta(1-a)=a(1-2\eta)+\eta$, which is the stated affine map; it is
sign-preserving around $\tfrac12$ for $\eta<\tfrac12$ and not invertible without $\eta$.
\end{proof}

\begin{proof}[Proof of Theorem~\ref{thm:rate-ub}]
Validity is Theorem~\ref{thm:cert} with the Maurer--Pontil deviation bound
$|\bar X-\Delta|\le\sqrt{2\widehat V\log(2/\alpha)/n}+7R\log(2/\alpha)/(3(n-1))$
\,w.p.\ $\ge1-\alpha$. For abstention: by Bernstein's inequality, w.p.\ $\ge1-\beta/2$,
$\bar X\ge\Delta-\sigma\sqrt{2\log(2/\beta)/n}-R\log(2/\beta)/(3n)$, and the empirical
variance satisfies $\widehat V\le 2\sigma^2+cR^2\log(2/\beta)/n$ on a further
$1-\beta/2$ event. On the intersection the radius is at most
$2\sigma\sqrt{2\log(2/\alpha)/n}+c'R\log(2/\alpha)/n$, so
$\bar X-\mathrm{rad}>0$ whenever $\Delta$ exceeds the displayed rate; the freeze side is
symmetric. No step uses $\sigma$.
\end{proof}

\begin{proof}[Proof of Theorem~\ref{thm:rate-lb}]
\emph{Dense branch.} Take $X=\pm\sigma$ with probabilities $q_\pm=\tfrac12(1\pm\kappa/\sigma)$
(mean $\pm\kappa$, variance $\le\sigma^2$). Then
$\mathrm{KL}(P_+^{\otimes n}\Vert P_-^{\otimes n})\le 8n\kappa^2/\sigma^2$ for
$\kappa\le\sigma/\sqrt2$. From an $(\alpha,\beta)$-reliable rule build the test
$\psi=\mathbf1[\textsc{adapt}]$; its two error probabilities are each at most
$\alpha+\beta$. Bretagnolle--Huber gives
$2(\alpha+\beta)\ge\tfrac12 e^{-8n\kappa^2/\sigma^2}$, i.e.\
$\kappa\ge\tfrac{\sigma}{2\sqrt2}\sqrt{\log\frac{1}{4(\alpha+\beta)}/n}$.
\emph{Spike branch (range-limited $1/n$ regime).} Let world$_-$ be the point mass at $-\kappa$ and world$_+$ place
mass $p$ at $R/2$ and $1-p$ at $-\kappa$ with $p(R/2+\kappa)=2\kappa$, so the means are
$\mp\kappa$. On the no-spike event --- probability $(1-p)^n$ --- the two worlds generate
identical samples, hence identical action laws. Reliability in world$_+$ forces commit
probability $\ge1-\alpha-\beta$ overall, so adapting in world$_-$ has probability at
least $(1-p)^n-(\alpha+\beta)$; validity caps this by $\alpha$. Therefore
$(1-p)^n\le2(\alpha+\beta)$, i.e.\ $n\log\frac{1}{1-p}\ge\log\frac{1}{2(\alpha+\beta)}$;
since $\log\frac{1}{1-p}\le \frac{p}{1-p}$ this gives $np\ge(1-p)\log\frac{1}{2(\alpha+\beta)}$, and
$\kappa=\tfrac{p}{2}(R/2+\kappa)\ge\tfrac{R}{4}(1-p)\,\log\frac{1}{2(\alpha+\beta)}/n$.
This branch certifies the separate range-limited $1/n$ regime; the headline $m^{-1/2}$
minimax rate is supplied by the dense branch above (matching the Maurer--Pontil upper
bound), the overall lower bound being the maximum of the two constructions.
\end{proof}
